% STATUS: written 2026-10-12
% Sources read in full: Kharchenko 1978 (pp. 155--156, the introduction and Section 1), Kharchenko 1991 (Section 2.4, pp. 112--116), Martindale 1969 (p. 576), Lanski 1988 (pp. 275--279).
\section{The theory of Kharchenko: differential identities and inner derivations}\label{sec:prime-semiprime:kharchenko}

Section~\ref{sec:prime-semiprime:quotient-rings} gave a field $C$ and a ring of
quotients $Q$. We ask which identities the values of derivations can satisfy in a
prime ring. Kharchenko asked this question in 1978 \cite[p.~155]{Kharchenko1978}.

By a \emph{differential identity} Kharchenko means an expression in which there
appear, besides unknowns and ring operations, elements of the ring and
derivations viewed as unary operations \cite[p.~155]{Kharchenko1978}.

Some of these identities hold in every ring. Kharchenko lists five forms of them
\cite[p.~155]{Kharchenko1978}. Three are the rule for a derivation of a product,
the formula of an inner derivation, and the formula for the commutator of two
derivations \cite[p.~155]{Kharchenko1978}. Two more are the rule for the $p$-th
power of a derivation in characteristic $p$, and the rules for coefficients that
are central \cite[p.~155]{Kharchenko1978}. Every generalized identity is also a
differential identity \cite[p.~155]{Kharchenko1978}.

We state the main result: nothing else occurs.

\begin{theorem}[{\cite[p.~155]{Kharchenko1978}}]\label{thm:prime:kharchenko-identities}
Let $R$ be a prime ring. Every multilinear differential identity of $R$, with
coefficients in the ring of quotients of $R$, is a consequence of identities of the
five forms above and of generalized identities satisfied in $R$.
\end{theorem}

Let a prime ring satisfy a differential identity that is not a consequence of the
five forms. Then it satisfies a nontrivial generalized identity
\cite[p.~155]{Kharchenko1978}. By the theorem of Martindale, its central closure
is a primitive ring with nonzero socle \cite[p.~155]{Kharchenko1978}. Its skew
field is finite dimensional over the center.

The book of Kharchenko gives the same result in the language of the previous
section. It shows that a differential identity of a prime ring $R$, written in
reduced form, yields a generalized identity on the ring of quotients $R_F$
\cite[Theorem~2.4.1, pp.~112--113]{Kharchenko1991}.

A derivation of $R$ is \emph{inner for the quotient ring} if its extension to $Q$
is an inner derivation of $Q$ \cite[p.~156]{Kharchenko1978}. The book proves a
lemma on such derivations \cite[Lemma~2.4.2, p.~113]{Kharchenko1991}.

\begin{theorem}[{\cite[Lemma~2.4.2, p.~113]{Kharchenko1991}}]\label{thm:prime:kharchenko-inner}
Let $R$ be a prime ring that satisfies a nontrivial generalized identity, and let
$C$ be its generalized centroid. Then every derivation of $R$ that is zero on
$C$ is inner for the ring of quotients.
\end{theorem}

Thus two derivations of $R$ that agree on $C$ differ by a derivation that is
inner for $Q$ \cite[p.~114]{Kharchenko1991}. So the derivations that $R$ induces
on $C$ show how independent the derivations of $R$ are
\cite[p.~114]{Kharchenko1991}.

In characteristic zero, every algebraic derivation is inner for the ring of
quotients \cite[p.~155]{Kharchenko1978}.
Section~\ref{sec:prime-semiprime:engel-nilpotent} returns to algebraic
derivations.

Kharchenko adds two results in his introduction \cite[p.~155]{Kharchenko1978}.
He shows that the outer derivations generate the universal enveloping algebra
of their Lie algebra. He also shows that every restricted Lie algebra is the
algebra of outer derivations of some prime ring.

Lanski used this method for the theorems of Posner \cite[p.~276]{Lanski1988}. He
shows that the derivations involved are inner. Then the ring satisfies a
generalized polynomial identity. The theorem of Martindale reduces the problem
to matrix rings over a field \cite[p.~276]{Lanski1988}.

Section~\ref{sec:prime-semiprime:engel-nilpotent} applies this reduction to Engel
conditions and nilpotent derivations.

A derivation can be constrained by a polynomial rather than a product. Lee and
Shiue take a polynomial that is not central-valued and two derivations $d$ and
$\delta$ whose action on its values differs by a central element, and prove that
one of the derivations is then central
\cite[Abstract]{LeeShiue1998Cocentralizing}.

